% ==============================================================================
% Chapter 2: Theoretical Foundation - Political Communication & Political Science
% ==============================================================================
\chapter{תשתית תיאורטית: תקשורת פוליטית, מדעי המדינה ולוחמת תודעה}
\label{ch:theory}

\section{מדוע מערכת AI לאומית חייבת להישען על תיאוריה פוליטית ותקשורתית?}
טעות תכנונית נפוצה בפיתוח מערכות לאיתור מידע כוזב היא התייחסות לדיסאינפורמציה כאל בעיית סיווג בינארית של משפטים ("אמת" מול "שקר"). מנקודת המבט של חקר התקשורת הפוליטית ומדעי המדינה, מסגור כזה מחמיץ את מהות התופעה: קמפיינים של השפעה זרה ומניפולציה פוליטית אינם מבוססים על עובדה בודדת המנותקת מהקשרה, אלא על שינוי מבני של תפיסת המציאות, קביעת סדר היום של הבוחרים, החרפת ניכור רגשי בין קבוצות, והרס התנאים המוקדמים לקיומו של שיח דמוקרטי משותף \cite{lazer2018, walker2018}.

פרק זה מציג את שבעת נדבכי התיאוריה האקדמית שעליהם בנויה ארכיטקטורת \en{CivicShield AI}. כל תיאוריה אינה מוצגת כרקע מופשט בלבד, אלא מתורגמת במישרין לאותות ממוחשבים (\en{Computational Signals}), מדדים כמותיים וסוכני ניתוח ייעודיים בתוך ועדת הסוכנים של המערכת.

\section{תיאוריות הליבה בתקשורת פוליטית}

\subsection{תיאוריית קביעת סדר היום והתערבות בין-תקשורתית (Agenda Setting)}
מאז מחקרם הקלאסי של מקומבס ושו (\en{McCombs \& Shaw, 1972}) בבחירות בצ'אפל היל, מקובל כי תפקידה המרכזי של התקשורת אינו בהכרח לקבוע לציבור \emph{מה} לחשוב, אלא \emph{על מה} לחשוב \cite{mccombs1972}. במרחב ההיברידי העכשווי, הדינמיקה עברה שתי תמורות מהותיות:
\begin{itemize}
    \item \textbf{סדר יום מדרגה שנייה (\en{Second-Level Attribute Agenda Setting}):} מיקוד תשומת הלב של הציבור לא רק בנושא עצמו (למשל, ניהול הקלפיות), אלא בהבלטת תכונות ספציפיות ("השערורייה", "קריסת המערכות") המטות את השיפוט האזרחי.
    \item \textbf{קביעת סדר יום בין-תקשורתית הפוכה (\en{Reverse Intermedia Agenda Setting}):} נחילי בוטים בקבוצות טלגרם ורשתות מיקרו-בלוגינג מייצרים גל מלאכותי של אזכורים תוך דקות ספורות; עיתונאים המנטרים מגמות רשת מאמצים את הנושא למהדורות המרכזיות מתוך חשש להחמיץ סקופ, ובכך מכשירים מבצע השפעה זר לכדי "שיח לגיטימי".
\end{itemize}
בתוך \en{CivicShield AI}, עיקרון זה ממומש באמצעות מעקב אחר גרדיאנט השיח הרב-ערוצי: מדידת זמן המעבר בין שיאים סינכרוניים של חשבונות חשודים ב-\en{Telegram/X} לבין הופעת הביטויים באתרי חדשות מרכזיים (\en{Mainstream Media Spillover}).

\subsection{פרדיגמת המסגור הפוליטי של אנטמן (Entman Framing Paradigm)}
רוברט אנטמן הגדיר מסגור כפעולה של בחירת היבטים מסוימים במציאות נתפסת והפיכתם לבולטים יותר בטקסט תקשורתי \cite{entman1993}. על פי אנטמן, כל מסגרת אפקטיבית ממלאת ארבע פונקציות מוגדרות: \textbf{(1) הגדרת הבעיה} (מה קורה); \textbf{(2) אבחון הסיבות} (מי אשם); \textbf{(3) שיפוט מוסרי} (האם המעשה לגיטימי או פסול); ו-\textbf{(4) המלצה לפתרון}.

מבצעי השפעה מתוחכמים נמנעים לעיתים משקרים עובדתיים קלים להפרכה, ומשקיעים במקום זאת ב\textbf{הזרקת מסגור עוין}. לדוגמה, תקלת חשמל מקומית במרכז מיון בוחרים תוצג עובדתית כנכונה (קרתה הפסקת חשמל בת 12 דקות), אך המסגרת שתיצוק לתוכה טענה של "חבלה מכוונת להעלמת קולות" משנה לחלוטין את האפקט הפוליטי. סוכן המסגור במערכת מפרק כל אשכול שיח חשוד לארבע רכיבי אנטמן באמצעות מודלי שפה ייעודיים ומזהה תאום סינטטי בהפצת מסגרות משולבות.

\begin{theorybox}[תרגום מודל אנטמן למערכת CivicShield AI]
במקום לבדוק "האם התרחש אירוע אמת?", סוכן המסגור (\en{Framing Agent}) משווה בין הממצא הפורנזי הטכני לבין תיוג הסיבתיות והשיפוט המוסרי המופצים במקביל במאות חשבונות. פער שיטתי ופתאומי בין הנתון העובדתי לבין עוצמת האשמת המוסדות מסומן כ\textbf{הסטת מסגור מאורגנת (\en{Coordinated Frame Injection})}.
\end{theorybox}

\subsection{ספירלת השתיקה, תיבות תהודה וקיטוב רגשי}
אליזבת נואלה-נוימן הראתה בתיאוריית \textbf{ספירלת השתיקה (\en{Spiral of Silence})} כי אנשים נמנעים מלהביע עמדה כאשר הם חשים שדעתם היא בדעת מיעוט, מתוך פחד מבידוד חברתי \cite{noelle1974}. המנגנון נשען על "איבר דמוי-סטטיסטי" (\en{Quasi-Statistical Organ}) שבעזרתו האזרח אומד את אקלים הדעות סביבו.

כאן נכנסת תופעת הרשתות החברתיות ונחילי הבוטים: הצפה מתואמת של אלפי תגובות קולניות ומאיימות בתגובה לידיעה פוליטית יוצרת \emph{אשליה אופטית} של רוב חברתי מדומה (\en{Illusory Majority}). כתוצאה מכך, אזרחים מתונים נסוגים מהשתתפות בדיון הפומבי — תהליך שמאיץ התכנסות לתוך \textbf{תיבות תהודה (\en{Echo Chambers})} \cite{sunstein2017} ומחריף את \textbf{הקיטוב הרגשי (\en{Affective Polarization})} — המצב שבו יריבים פוליטיים אינם נתפסים עוד כבעלי דעה שונה, אלא כאויב מוסרי שיש לשנוא ולפסול \cite{iyengar2019}.

\subsection{תיאוריית החיסון התודעתי (Inoculation Theory / Prebunking)}
הפריצה הקוגניטיבית של מקגווייר משנת 1961 \cite{mcguire1961}, שעובדה בעשור האחרון לסביבת הרשת בידי רוזנבק, ואן דר לינדן ולוונדובסקי \cite{roozenbeek2022}, משמשת כבסיס ההתערבות המרכזי של \en{CivicShield AI}. הרעיון בנוי משני מרכיבים:
\begin{enumerate}
    \item \textbf{אזהרה מוקדמת (\en{Forewarning}):} אזהרה ניטרלית לציבור כי צפויה פנייה רגשית מניפולטיבית ("בתקופת בחירות נפוצים ניסיונות לזייף התבטאויות מוקלטות ולגרום לפאניקה מלאכותית").
    \item \textbf{הפרכה מונעת של הטכניקה (\en{Refutational Preemption}):} חשיפת הטכניקה הרטורית ("הפחדה באמצעות הקשר חתוך", "חיקוי קול בסינתזת AI") על גבי דוגמה מוחלשת, המעוררת את ההגנות הקוגניטיביות של האזרח בעת מפגש עתידי עם הנרטיב האמיתי.
\end{enumerate}

מחקרי מעבדה ושדה רחבי היקף הוכיחו כי חיסון מוקדם משפר את יכולת האזרחים לזהות מניפולציות ב-20\% עד 35\% ומקטין שיעורי שיתוף של תוכן מסית, ללא כל תלות בזהות המפלגתית של הנבדקים \cite{roozenbeek2022}.

\section{נדבכים במדעי המדינה ובמחשבה החוקתית}

\subsection{המרחב הציבורי של הברמאס והאפיסטמולוגיה הדמוקרטית}
יורגן הברמאס תיאר את \textbf{המרחב הציבורי (\en{Public Sphere})} כתחום החיים החברתיים שבו אזרחים מתכנסים כשווים לשיח רציונלי-ביקורתי המוליד דעת קהל לגיטימית \cite{habermas1989}. קיומו של המרחב הציבורי מתאפשר רק בהינתן תקשורת שאינה מעוותת באורח שיטתי בידי שחקני כוח סמויים. כאשר אלגוריתמי רשת משדכים תיבות תהודה וכאשר שחקנים זרים מזריקים פרוקסי אוטומטי לתוך הדיון, המרחב מאבד את הדו-כיווניות הרציונלית והופך לתהליך של "פיאודליזציה מחודשת" — שבו הכוח הכספי או הטכנולוגי כובש את כושר השיפוט האזרחי.

\subsection{עוצמה חדה (Sharp Power) בעידן הדיגיטלי}
ג'וזף ניי הבחין בין \textbf{עוצמה קשה (\en{Hard Power})} — כפייה צבאית וכלכלית — לבין \textbf{עוצמה רכה (\en{Soft Power})} הנשענת על משיכה תרבותית וערכית \cite{nye2004}. לעומתם, כריסטופר ווקר וג'סיקה לודוויג טבעו ב-2018 את המושג \textbf{עוצמה חדה (\en{Sharp Power})} לתיאור מדיניות של משטרים אוטוריטריים המנצלים את הפתיחות האינפורמטיבית של דמוקרטיות כדי לנקב, לחדור ולעוות את המערכת הפוליטית מבפנים \cite{walker2018}. עוצמה חדה אינה מנסה "לשכנע" באטרקטיביות של המדינה התוקפת, אלא \emph{לשבש את האמון} של האזרחים במוסדותיהם שלהם ("הכול מושחת, אין טעם להצביע").

\subsection{דמוקרטיה מתגוננת (Militant / Defensive Democracy)}
המשפטן והמדען המדינה קרל לוינשטיין טבע ב-1937 את דוקטרינת "הדמוקרטיה הנאבקת" (\en{Militant Democracy}) כלקח מקריסת רפובליקת ויימאר: דמוקרטיה אינה חייבת להפוך עצמה למצע חסר הגנה להרס עצמי של עקרונותיה \cite{loewenstein1937}. בישראל פיתח השופט אהרן ברק מסורת פסיקתית מאזנת של "דמוקרטיה מתגוננת" המכירה בזכות המדינה להגן על עצם קיומה הדמוקרטי וטוהר הבחירות, אך תוך \textbf{ריסון מידתי קפדני} ושמירה עליונה על חופש הביטוי הפוליטי \cite{barak1988}.

\begin{riskbox}[הסכנה החוקתית: מ"דמוקרטיה מתגוננת" למדינת צנזורה אלגוריתמית]
אם סמכות הזיהוי של \en{AI} תקושר אוטומטית לסמכות ההסרה של התבטאות פוליטית, הטכנולוגיה עלולה לייצר "אפקט מצנן" חמור על מתנגדי שלטון ומיעוטים. לפיכך, המערכת נאמנה למשוואה החוקתית של ברק: ניטור אוטומטי וחיסון ציבורי מותרים; סנקציה של הגבלת תוכן מחייבת לעולם חתימה שיפוטית גלויה ובת-ערעור (פרק~\ref{ch:ethics}).
\end{riskbox}

\section{מיפוי התיאוריות לארכיטקטורת CivicShield AI}
השלב המסכם של החיבור התיאורטי מוצג באיור~\ref{fig:public_sphere_comparison}, המדגים מבנית את ההבדל בין מרחב ציבורי תקין לבין מרחב הנתון למתקפת \en{FIMI} משולבת, ולאחריו טבלה~\ref{tab:theory_to_code} המתרגמת כל תיאוריה ישירות לסוכן \en{AI} ולאלגוריתם מדידה.

\begin{figure}[htbp]
\begin{tikzhebrew}
\begin{tikzpicture}[
    node distance=0.6cm and 0.6cm,
    cleanbox/.style={rectangle, rounded corners=4pt, draw=CivicTeal, fill=CivicTealBg, line width=1pt, minimum width=5.6cm, minimum height=1.15cm, align=center, font=\footnotesize},
    poisonbox/.style={rectangle, rounded corners=4pt, draw=CivicRed, fill=CivicRedBg, line width=1pt, minimum width=5.6cm, minimum height=1.15cm, align=center, font=\footnotesize},
    hdr/.style={rectangle, rounded corners=3pt, minimum width=6.2cm, minimum height=0.75cm, align=center, font=\bfseries\small, text=white},
    arrteal/.style={-{Stealth[length=2.6mm]}, line width=1pt, color=CivicTeal},
    arrred/.style={-{Stealth[length=2.6mm]}, line width=1pt, color=CivicRed}
]
    % Left Side: Clean Habermasian Sphere (with CivicShield Prebunking)
    \node[hdr, fill=CivicTeal] (lhdr) at (-3.7, 4.2) {\he{א. מרחב ציבורי חסין (דליברטיבי)}};
    \node[cleanbox] (l1) at (-3.7, 2.8) {\textbf{\he{מגוון מקורות ומסגור מאוזן}}\\[2pt]{\scriptsize\he{דיווח עיתונאי שקוף ואחריותיות מוסדית}}};
    \node[cleanbox] (l2) at (-3.7, 1.2) {\textbf{\he{חיסון קוגניטיבי (\en{Prebunking})}}\\[2pt]{\scriptsize\he{חשיפה מוקדמת לטכניקות מניפולציה והנדסה}}};
    \node[cleanbox] (l3) at (-3.7, -0.4) {\textbf{\he{דיאלוג חוצה-קבוצות מאומת}}\\[2pt]{\scriptsize\he{סיכול בוטים ושמירת חופש ביטוי לאזרחים}}};
    \node[cleanbox, fill=CivicNavyBg, draw=CivicNavy] (l4) at (-3.7, -2.0) {\textbf{\he{תוצאה: לגיטימציה וטוהר בחירות}}\\[2pt]{\scriptsize\he{שימור הסמכות האפיסטמית והכרעה חופשית}}};

    \draw[arrteal] (l1) -- (l2);
    \draw[arrteal] (l2) -- (l3);
    \draw[arrteal] (l3) -- (l4);

    % Center Divider
    \draw[dashed, line width=1.1pt, color=CivicSlate] (0, 4.6) -- (0, -2.8);
    \node[rectangle, fill=white, draw=CivicSlate, rounded corners=2pt, font=\scriptsize\bfseries] at (0, 0.9) {\en{VS}};

    % Right Side: Poisoned Sphere under FIMI
    \node[hdr, fill=CivicRed] (rhdr) at (3.7, 4.2) {\he{ב. מרחב מורעל תחת מתקפת \en{FIMI}}};
    \node[poisonbox] (r1) at (3.7, 2.8) {\textbf{\he{חטיפת סדר יום (\en{Agenda Hijacking})}}\\[2pt]{\scriptsize\he{נחילי בוטים מתואמים והזרקת מסגרות עוינות}}};
    \node[poisonbox] (r2) at (3.7, 1.2) {\textbf{\he{\en{Deepfake} והתגברות \en{Sharp Power}}}\\[2pt]{\scriptsize\he{ייצור משבר אמון פתאומי ב-72 השעות שלפני הקלפי}}};
    \node[poisonbox] (r3) at (3.7, -0.4) {\textbf{\he{תיבות תהודה וספירלת השתיקה}}\\[2pt]{\scriptsize\he{רוב מדומה מרתיע מתונים ומעמיק קיטוב רגשי}}};
    \node[poisonbox, fill=CivicRed!18!white] (r4) at (3.7, -2.0) {\textbf{\he{תוצאה: שחיקת הלגיטימציה הדמוקרטית}}\\[2pt]{\scriptsize\he{אובדן מציאות עובדתית משותפת וניכור מוסדי}}};

    \draw[arrred] (r1) -- (r2);
    \draw[arrred] (r2) -- (r3);
    \draw[arrred] (r3) -- (r4);
\end{tikzpicture}
\end{tikzhebrew}
\caption{השוואה מבנית בין מרחב ציבורי דמוקרטי חסין לבין ספרה ציבורית תחת מתקפת FIMI}
\label{fig:public_sphere_comparison}
\end{figure}

{\renewcommand{\arraystretch}{1.35}
\begin{table}[htbp]
\centering
\caption{טבלת תרגום תיאוריות בתקשורת ובמדעי המדינה לרכיבים אלגוריתמיים}
\label{tab:theory_to_code}
\begin{tabularx}{\textwidth}{|p{2.9cm}|p{2.7cm}|p{3.0cm}|X|}
\hline
\textbf{\textcolor{CivicNavy}{תיאוריה וחוקרים}} & \textbf{\textcolor{CivicNavy}{מנגנון המניפולציה}} & \textbf{\textcolor{CivicNavy}{רכיב ב-\en{CivicShield}}} & \textbf{\textcolor{CivicNavy}{מדד כימות ויישום חישובי}} \\ \hline
קביעת סדר יום \cite{mccombs1972} & חטיפת סדר יום בין-תקשורתית & \en{Network \& Trend Agent} & מדידת פער זמן זליגה מטלגרם לחדשות וגרדיאנט נרטיב (\en{Shift Rate}). \\ \hline
מסגור פוליטי (אנטמן) \cite{entman1993} & הזרקת אשם ופתרון מקטב & \en{Framing Analyzer Agent} & סיווג \en{NLP} ל-4 פונקציות אנטמן ואיתור אשכולות העתקה סמנטית. \\ \hline
ספירלת השתיקה \cite{noelle1974} & אפקט רוב מזויף והשתקת מיעוט & \en{CIB / Bot Detector} & יחס קול/משתתף אנושי, מתאם סינכרוני (\en{$\Delta t < 2s$}) וחריגת רשת. \\ \hline
קיטוב רגשי ותיבות תהודה \cite{sunstein2017, iyengar2019} & פיצול רשתי להקצנה רגשית & \en{Simulator (H-K Engine)} & מודל \en{Hegselmann-Krause}, מרחק דו-אופני (\en{Bimodality}) ואנטרופיית דיאלוג. \\ \hline
חיסון תודעתי (\en{Prebunking}) \cite{roozenbeek2022} & פגיעות טרום-חשיפה & \en{Prebunking Composer} & מודל \en{SIRI ODEs}, חישוב סף התפרצות \en{$R_0$} ושיעור חיסון קריטי \en{$p_c$}. \\ \hline
עוצמה חדה (\en{Sharp Power}) \cite{walker2018} & שיוך זר וחבלה במוסדות & \en{Forensic Attribution Agent} & איחוד בייסיאני של טביעות אצבע לשוניות, אודיו וזמני פעילות גאוגרפיים. \\ \hline
דמוקרטיה מתגוננת \cite{loewenstein1937, barak1988} & סכנת צנזורה יתר אלגוריתמית & \en{Constitutional Guardrail} & שער חתימה משולשת, אכיפת \en{$\epsilon$-DP} ואיסור הסרה ללא שופט. \\ \hline
\end{tabularx}
\end{table}}

\section{סיכום הפרק והיגזר התכנוני}
המעבר מתיאוריית התקשורת הקלאסית אל ארכיטקטורת הבינה המלאכותית מעניק למערכת \en{CivicShield AI} יתרון כפול: היא נמנעת מניסיון עקר להפוך ל"מיניסטריון אמת", ובמקום זאת מכוונת בדיוק אל הכשלים המבניים שאותם הוכיחה הספרות האקדמית — סינכרוניזציה מלאכותית, שיבוש סדר יום, הזרקת מסגור וניצול דינמיקה של ספירלת שתיקה. פרק~\ref{ch:requirements} יתרגם עקרונות אלה למפרט דרישות פונקציונלי מלא ותרחישי ייחוס אופרטיביים.
